\section{Assessment Scope}

This report documents a Data Management project focused on restaurant review
platforms in the Milan area. The current report is a pre-integration baseline:
it assesses Google Places, Tripadvisor, and TheFork data after acquisition and
before entity matching. The project can later integrate these sources to analyse
cross-platform rating consistency on a unified restaurant table.

\section{Data Sources}

The current acquisition layer contains three raw source datasets:
Google Places as geographic seed and reference source, Tripadvisor as a
general-purpose review platform, and TheFork as a restaurant-specific platform.
All quality metrics in this report are computed from raw acquisition outputs.
No entity matching or post-integration metric is included in this baseline.

\section{Pre-Integration Quality Metrics}

The selected data-quality dimensions are completeness, validity/consistency,
uniqueness, timeliness, and reliability. Completeness measures missing values;
validity checks source-specific field formats on present values; uniqueness
checks stable source identifiers; timeliness measures source refreshability from
collection duration; reliability flags ratings based on sparse review evidence.
All percentage component scores are oriented so that 100\% is better. This is a
raw pre-cleaning baseline and should be rerun after cleaning/enrichment to
measure improvement with the same formulas.

\subsection{Metric Definitions}

\preintegrationtable{metric_definitions.tex}

\subsection{Cross-Source Summary}

\preintegrationtable{source_summary.tex}

\subsection{Weighted Quality Score}

The weighted score favours dimensions required for integration and analysis:
critical completeness, validity, spatial readiness, uniqueness, timeliness, and
review-count reliability. This avoids treating sparse optional metadata as equal
to identifiers, ratings, review counts, and coordinates. Average completeness is
reported separately, while critical completeness is the completeness component
used in the weighted score.

\preintegrationtable{source_quality_scores.tex}

\subsection{Pre-Integration Visual Diagnostics}

The following generated visuals summarize source readiness before matching:
weighted score bars, component-level diagnostics, core-field coverage, and the
most frequent anomaly classes.

\subsubsection*{Weighted quality score}
\preintegrationtable{visual_quality_scores.tex}

\subsubsection*{Score components}
\preintegrationtable{visual_score_components.tex}

\subsubsection*{Core coverage heatmap}
\preintegrationtable{visual_coverage_heatmap.tex}

\subsubsection*{Anomaly profile}
\preintegrationtable{visual_anomaly_profile.tex}

\subsection{Comparative Findings}

\preintegrationtable{source_comparison.tex}

\subsection{Quality Improvement Actions}

\preintegrationtable{improvement_actions.tex}

\section{Source-Specific Assessment}

\preintegrationtable{google_places_detail.tex}

\preintegrationtable{tripadvisor_detail.tex}

\preintegrationtable{thefork_detail.tex}

\section{Complete Field Coverage by Source}

\preintegrationtable{google_places_field_coverage.tex}

\preintegrationtable{tripadvisor_field_coverage.tex}

\preintegrationtable{thefork_field_coverage.tex}

\section{Next Steps}

After geospatial integration and entity matching, this report should be extended
with match-rate metrics, unmatched and ambiguous record counts, and
cross-platform rating-disagreement metrics on the unified restaurant table. In
particular, Tripadvisor latitude and longitude enrichment is required before the
geospatial integration step can be evaluated fairly.

\section{Reproducibility}

The report is regenerated by running \texttt{report/build\_report.ps1} from
PowerShell. The script profiles the three raw datasets, refreshes the Markdown
and LaTeX tables, and then compiles \texttt{report/main.pdf}.
